\section{Proof of the all-chain transfer theorem}
\label{app:transfer-proof}

This appendix proves \cref{thm:transfer} from the finite local hypotheses. All constants below depend only on the fixed palette.

\subsection{Global decomposition and interface energy}

Write the target chain space as the orthogonal coordinate sum
\[
C_d(X_A)=C_d(R)\oplus\bigoplus_{j\in A}C_{d,j}^{\rm priv}.
\]
For a unit chain \(x\), write
\[
x=\omega_0+\sum_j\omega_j,
\qquad
x_j=\omega_0+\omega_j,
\qquad
U=\sum_j\|\omega_j\|^2,
\qquad
e=\langle x,\Delta_Ax\rangle.
\]
Let \(a=\partial_R\omega_0\), \(y_j=L_j\omega_j\), \(y=\sum_jy_j\), and \(r=a+y\). The vector \(r\) is the shared-register component of the global down-boundary. Since
\(
a+y_j=r-(y-y_j)
\),
\[
\sum_j\|a+y_j\|^2
\le2t_A\|r\|^2+2\sum_j\|y-y_j\|^2.
\]
The elementary identity
\[
\sum_j\|y-y_j\|^2
=(t_A-2)\|y\|^2+\sum_j\|y_j\|^2
\]
also holds for \(t_A=1\). Put
\[
\rho=\lambda^{-2}\|[L_1\ \cdots\ L_{t_A}]\|^2
\le N_*t_A,
\qquad
\chi=(t_A-2)_+\rho+N_*\le C t_A^2.
\]
Then~\cref{eq:interface} gives
\[
\sum_j\|y-y_j\|^2\le\chi\lambda^2U.
\]
All nonregister down-boundary coordinates and all up-boundary coordinates are private and partition by gadget. We obtain
\begin{equation}
\label{eq:appendix-local-energy}
\sum_j\|D_j(\lambda)x_j\|^2
\le2t_Ae+2\chi\lambda^2U.
\end{equation}

\subsection{Zero-weight and private-sector coercivity}

Let \(A_{j,0}=I-P_V-Q_j\) in the padded local chain space. The positive singular-value floor of \(D_{j,0}\) on \(\ran A_{j,0}\), together with
\(
D_j(\lambda)=D_{j,0}+\lambda D_{j,1}
\), gives
\begin{align}
S_A:=\sum_j\|A_{j,0}x_j\|^2
&\le C\sum_j\|D_{j,0}x_j\|^2\notag\\
&\le C\left(
t_Ae+\chi\lambda^2U+
\lambda^2\sum_j\|x_j\|^2
\right)\notag\\
&\le C\left(t_Ae+\chi\lambda^2U+t_A\lambda^2\right).
\label{eq:appendix-SA}
\end{align}
The final inequality uses
\(
\sum_j\|x_j\|^2
=t_A\|\omega_0\|^2+U\le t_A
\).

Because \(T_jP_V=0\),
\[
T_jQ_jx_j=T_jx_j-T_jA_{j,0}x_j.
\]
The lower bound on \(T_j\) restricted to \(Q_j\), the estimate \(\|T_j\|\le C\lambda\), and orthogonality of the retained private output blocks imply
\begin{equation}
\label{eq:appendix-SQ}
S_Q:=\sum_j\|Q_jx_j\|^2
\le C\left(\frac{e}{\lambda^2}+S_A\right).
\end{equation}
The orthogonality here concerns the retained \(T_j\) outputs. The shared \(L_j\) outputs were handled by their horizontal operator norm and need not be orthogonal.

Let \(N=\|(I-P_V)x\|^2\). Summing the orthogonal decompositions \(P_V\oplus Q_j\oplus\ran A_{j,0}\) gives
\[
N\le U+t_A\|(I-P_V)\omega_0\|^2=S_A+S_Q.
\]
Since \(U\le N\), substituting~\cref{eq:appendix-SA,eq:appendix-SQ} and using \(\chi\le Ct_A^2\) yields
\[
N\le C\left[t_A\lambda^2+(t_A+\lambda^{-2})e\right]
+Ct_A^2\lambda^2N.
\]
Choosing \(\lambda\le c_0/t_{\max}\) absorbs the last term and proves~\cref{eq:geometric-leakage}.

\subsection{Logical coercivity}

For each gadget, exact filling gives \(\ran\Pi_j\subseteq B_d(Y_j)\). The up-Laplacian has kernel \(B_d(Y_j)^\perp\), and~\cref{eq:local-gap} lower-bounds its positive spectrum. Therefore, after the common embedding into \(X_A\),
\[
\Delta_j^\uparrow\succeq c\lambda^\kappa\Pi_j.
\]
The top-dimensional register has no up-columns, and every degree-\((d+1)\) simplex with a private vertex belongs to exactly one gadget. Hence the global up-Laplacian is the exact sum of the embedded local up-Laplacians, proving~\cref{eq:logical-domination}.

Since \(K_A\subseteq V\), the squared distance decomposes as
\[
\|(I-P_{K_A})x\|^2
=N+\langle x,(P_V-P_{K_A})x\rangle.
\]
\Cref{eq:logical-gap,eq:logical-domination} give
\[
\langle x,(P_V-P_{K_A})x\rangle
\le\frac{e}{cg_A\lambda^\kappa}.
\]
Together with~\cref{eq:geometric-leakage},
\[
\|(I-P_{K_A})x\|^2
\le C\left[t_A\lambda^2+(t_A+\lambda^{-2})e
+\frac{e}{g_A\lambda^\kappa}\right].
\]
For \(\lambda\le c_0/t_{\max}\), \(g_A\le1\), and \(\kappa\ge2\), the last energy term dominates the middle energy terms up to a palette constant. This proves~\cref{eq:transfer-bound}.

\subsection{Exact multiplicity and the whole-positive-spectrum gap}

Let \(\mathcal L_A(E_A)\) be the span of all eigenvectors of \(\Delta_A\) with eigenvalue below \(E_A\). Every unit \(x\in\mathcal L_A(E_A)\) obeys \(e<E_A\). Choosing the constants in~\cref{eq:lambda-E-choice} makes the right-hand side of~\cref{eq:transfer-bound} strictly less than \(\eta^2<1\). Therefore \(P_{K_A}x\neq0\) for every nonzero \(x\in\mathcal L_A(E_A)\), so \(P_{K_A}|_{\mathcal L_A(E_A)}\) is injective and
\begin{equation}
\label{eq:low-dimension}
\dim\mathcal L_A(E_A)\le\dim K_A.
\end{equation}

Every \(v\in V\) remains a cycle after attaching the gadgets. By \cref{lem:independent-filling}, the map \(V/W_A\to H_d(X_A)\) is injective. Its domain has dimension \(\dim K_A\), so Hodge theory supplies at least \(\dim K_A\) zero eigenvectors of \(\Delta_A\). These zero vectors belong to \(\mathcal L_A(E_A)\). Combining this lower bound with~\cref{eq:low-dimension} forces equality throughout: the geometric kernel has dimension \(\dim K_A\), and no positive eigenvalue lies below \(E_A\).

If \(K_A=0\), any unit vector in \(\mathcal L_A(E_A)\) would have distance one from \(K_A\), contradicting the strict \(\eta<1\) bound. Thus the same argument proves \(\Delta_A\succeq E_AI\) in the zero-kernel case. No additional global relative-acyclicity assumption is required.
